%==========================================================================
\section{Test plan}
\begin{center}\small
\begin{tabularx}{\linewidth}{@{}>{\raggedright}p{3.4cm} >{\raggedright}X X@{}}
\toprule
\textbf{Test Case Name} & \textbf{Description} & \textbf{Expected Outcome} \\
\midrule
1. Subclass \& Parent Introspection &
  Stores a subclass instance in a parent-typed variable, calls
  \code{reflect()} on it, and prints \code{class\_name()} and
  \code{parent\_name()}. &
  Prints the subclass name, then the parent name. Exit code 0. \\
\midrule
2. Inherited Field \& Overridden Method Enumeration &
  Queries ReflectInfo field and method counts on a class with inherited
  members and overridden methods. &
  Fields are listed parent-first; each override appears once; arities and
  signatures print as expected. Exit code 0. \\
\midrule
3. Scalar Field Read/Write Round-trip &
  Performs setfield and getfield on int and bool fields using dynamic string
  names. &
  Values mutate correctly without corrupting adjacent memory offsets. Exit
  code 0. \\
\midrule
4. Reference Field Mutation \& GC Aliasing &
  Uses setfield to set a class-typed field to a subclass instance and then to
  null, triggering GC in between. &
  Reference identity preserved; object kept alive during GC cycle; null
  stored safely. Exit code 0. \\
\midrule
5. Multi-Argument Method Invocation &
  Calls methods by name with several arguments (including a subclass value
  passed to a parent-typed parameter), in both the value and void forms. &
  The overriding method runs, arguments arrive correctly, and the expected
  result is returned. Exit code 0. \\
\midrule
6. Reflective Construction via new\_instance &
  Constructs a subclass by name with \code{new\_instance(Parent, "Child", \ldots)}
  using constructor arguments. &
  The object is allocated and initialized (including String fields defaulting
  to \code{""}), and \code{class\_name()} reports Child. Exit code 0. \\
\midrule
7. Static Semantics Negatives &
  Section 4 ill-formed program: \code{setfield(int, b, "count", true)}. &
  Rejected at compile time with \code{E\_REFLECT\_VALUE\_TYPE}. \\
\midrule
8. Dynamic Validation Failure Sequence &
  Separate runs with a missing member name, a wrong requested type, a wrong
  argument count, a new\_instance or allocate of a class that is not a
  subclass of T, and a setfield of a plain Serial into a Node-typed field. &
  Each aborts with exit code 140 and the matching \code{lo\_reflect: <reason>};
  no field is written and no body runs. \\
\midrule
9. Null Receiver Side-effects &
  Calls invoke on a null receiver, with arguments that have visible side
  effects (e.g.\ printing); a separate run calls \code{reflect(null)}. &
  The argument side effects occur left to right, then the program aborts with
  exit code 140 (null receiver); \code{reflect(null)} aborts the same way. \\
\midrule
10. Supertype Field Access &
  With field next declared Node (Node extends Serial), prints
  \code{getfield(Serial, o, "next") = n} and calls
  \code{setfield(Serial, o, "next", m)} with a Node m. &
  Prints true; the setfield succeeds. Exit code 0. \\
\midrule
11. Reflective Allocation via allocate &
  In a loop that forces collections, calls \code{allocate(Parent, "Child")},
  where Child's constructor prints a line, and prints \code{class\_name()}
  and each field. &
  Prints Child, then 0, false, \code{""}, and true for \code{(field = null)};
  the constructor's line never appears. Exit code 0. \\
\midrule
12. Serializer Round-trip &
  Serializes a small Serial graph with a shared node and a cycle,
  deserializes it, and prints the rebuilt graph's fields. &
  Output matches the original's printout; the shared node appears once. Exit
  code 0. \\
\bottomrule
\end{tabularx}
\end{center}

